% Spacings between the chapter top and to the text
%-----------------------------------------------------------------------------
\newcommand{\chaptertopspacing}{10}
\newcommand{\chaptertotextspacing}{-10}
\newcommand{\chaptertopspacingSTAR}{10}
\newcommand{\chaptertotextspacingSTAR}{-10}

\iffinal{}{  % Library format requires numbered footnotes
\renewcommand{\thefootnote}{$\dagger$}	% Daggers for footnote markers style
}

% Convenience functions
% -----------------------------------------------------------------------------
\newcommand{\B}[1]{ {\bf #1} }      		% Bold text for matrices
\newcommand{\Z}{\mathcal{Z}}				% Partition function
\newcommand{\HAM}{\mathcal{H}}		% Hamiltonian
\newcommand{\qp}[0]{(\B{q},\B{p})}  % Cannonical variables (q,p) 
\newcommand{\pfrac}[2]{\frac{\partial #1}{\partial #2}} 	% partial 1 /partial 2 
%\newcommand{\fr}{2}{\frac{#1}{#2}}   % frac 1/2
\newcommand{\ket}[1]  {\left | \, #1 \right \rangle}			% bra-ket notation
\newcommand{\abs}[1]{\left\vert #1 \right\vert}		% vert bars for averages
\newcommand{\norm}[1]{\left\Vert #1 \right\Vert}	% taller vert bars for the norm
\newcommand{\evalat}[1]{\left. #1 \right \vert}	% ex. evaluating the int. at its limits
\newcommand{\set}[1]{\left\{ #1 \right\}}					% squigle brackets for sets
\newcommand{\avg}[1]{\left< #1 \right>}					% angle brackets for averages
\newcommand{\paren}[1]{\left( #1 \right)}					% grows parentheses to the correct size
\newcommand{\brackets}[1]{\left[ #1 \right]}			% grows square brackets
\newcommand{\braces}[1]{\left \{ #1 \right \}}
\newcommand{\piecewisebrace}[1]{\left \{ #1 \right .} % Useful for piecewise functions

\newcommand{\ie}{i.e.\ }
\newcommand{\eg}{e.g.\ }
\newcommand{\cf}{c.f.\ }
\newcommand{\chem}[1]{\ensuremath{\mathrm{#1}}} % Chemical formula

\DeclareMathOperator*{\minor}{minor\,}
\DeclareMathOperator*{\sgn}{\mathbf{sgn}\,}
\DeclareMathOperator*{\rootof}{RootOf\,}


\newcommand{\figurewidthSINGLE}{0.8\textwidth}
\newcommand{\figurewidthDOUBLE}{0.45\textwidth}
\newcommand{\figurewidthTRIPLE}{0.35\textwidth}

\newcommand{\CHI}[2]{\chi_{#1 #2}} % Cluster Counting function

% ----------------------------------------------------------------------------

\newcommand{\civ}{\ion{C}{4}}
\newcommand{\athree}{$\alpha_{\rm 3D}$}
\newcommand{\aone}{$\alpha_{\rm 1D}$}
\newcommand{\atwo}{$\alpha_{\rm 2D}$}
\newcommand{\luvaox}{$L_{\rm }$--$\alpha_{\rm ox}$}
\newcommand{\lbol}{$L_{\rm bol}$}
\newcommand{\luv}{$L_{\rm uv}$}
\newcommand{\lx}{$L_{\rm x}$}
\newcommand{\llam}{$L_{\lambda}$}
\newcommand{\ltwofive}{$L_{\rm 2500\,\AA}$}
\newcommand{\lfiveone}{$L_{\rm 5100\,\AA}$}
\newcommand{\bctwofive}{BC$_{\rm 2500\,\AA}$}
\newcommand{\bcfiveone}{BC$_{\rm 5100\,\AA}$}
\newcommand{\onemum}{1\,$\mu$m}
\newcommand{\tenmum}{10\,$\mu$m}
\newcommand{\thirtymum}{30\,$\mu$m}
\newcommand{\twokev}{2\,keV}
\newcommand{\tenkev}{10\,keV}
\newcommand{\hundredev}{0.1\,keV}
\newcommand{\twohundredev}{0.2\,keV}
\newcommand{\ebv}{{\rm E(B-V)}}
\newcommand{\dal}{\Delta(\alpha_{\lambda})}
\newcommand{\dgi}{\Delta(g-i)}
\newcommand{\dug}{\Delta(u-g)}
\newcommand{\diz}{\Delta(i-z)}
\newcommand{\mrel}{\Delta m}
\newcommand{\mbh}{M_{\rm BH}}
\newcommand{\lmbh}{\log{(M_{\rm BH})}}
%\newcommand{\lf}{L_{\rm frac}}
\newcommand{\lf}{L/L_{\rm Edd}}

% ----------------------------------------------------------------------------
% Override the color of links
%\newif\if@colorref
\hypersetup{colorlinks}
%\if@colorref
  \hypersetup{linkcolor=blue}
  \hypersetup{anchorcolor=black}
  \hypersetup{citecolor=blue}
  \hypersetup{filecolor=black}
  \hypersetup{menucolor=black}
  \hypersetup{runcolor=black}
  \hypersetup{urlcolor=blue}
\fi

% Partition block notation
% -----------------------------------------------------------------------------
\newcommand*\pblock[3]{
	\tikz[baseline=(char.base)]{
		\node[ellipse,draw,inner sep=2pt] (char) {$#1$};
	} ^{#2}_{#3}
}

\newcommand{\Go}{\text{G}\bar{\text{o}}}
\newcommand{\GO}{\textbf{G}}
\newcommand{\CHIX}{\mathbf{X}}
\newcommand{\KP}{k_+}
\newcommand{\KM}{k_-}
\newcommand{\KPMAX}{k_+^*}

% -----------------------------------------------------------------------------
% Add the CC logo to the copyright page and spruce it up a bit
\renewcommand\copyrighttextCCBYSA{
\\[2 em]
\begin{quote}
\begin{center}
\includegraphics[width=.3\textwidth]{pictures/cc_logo.pdf} \\[2 em]
\end{center} 
    This work is licensed under the terms of the Creative Commons
    Attribution-ShareAlike 4.0 International license.  The license is
    available at \\
    \url{http://creativecommons.org/licenses/by-sa/4.0/}.
\end{quote}
}

% -----------------------------------------------------------------------------
% Environment for temporarily removing the page footer
\newcommand{\suppressfooter}
{%
\fancyfoot[RE,LO,LE,RO]{}
\renewcommand{\footrulewidth}{0.0pt}
}%

\newcommand{\restorefooter}
{%
\fancyfoot[RE,LO]{\scshape\leftmark}
\fancyfoot[LE,RO]{\scshape\rightmark}
\renewcommand{\footrulewidth}{.4pt}
}%

\newcommand{\rightfooter}
{%
%\fancyfoot[RE,LO]{\scshape\leftmark}
\fancyfoot[LE,RO]{\scshape\rightmark}
\renewcommand{\footrulewidth}{.4pt}
}%

% -----------------------------------------------------------------------------
% Simple command for inline figures (does not count towards a normal figure)

\newcommand{\inlinespacing}{1em}
\newcommand{\inlineframebox}[2]{
\vspace{1em}
\newline
\fbox{
\begin{minipage}[]{.95 \linewidth}
\begin{center}
#1 \\
#2
\end{center}
\end{minipage}
}
\vspace{1em}
\newline
}

\newcommand{\inlinembox}[2]{
\vspace{1em}
\newline
\begin{minipage}[]{.95 \linewidth}
\begin{center}
#1 \\
#2
\end{center}
\end{minipage}
\vspace{1em}
\newline
}


% -----------------------------------------------------------------------------
% Custom styling for listing the source code with syntax highlighting

\newcommand{\SyntaxHighLighting}{
% the following is needed for syntax highlighting
\usepackage{xcolor}
\usepackage{textcomp}

\lstdefinelanguage{Scheme}{
  morekeywords=[1]{define, define-syntax, define-param, define-macro, lambda, define-stream, stream-lambda},
  morekeywords=[2]{begin, call-with-current-continuation, call/cc,
    at-beginning, after-sources,
    call-with-input-file, call-with-output-file, case, cond,
    do, else, for-each, if,
    let*, let, let-syntax, letrec, letrec-syntax,
    set!, set-param!, run-sources+, run-until,
    let-values, let*-values,
    and, or, not, delay, force,
    quasiquote, quote, unquote, unquote-splicing,
    map, fold, syntax, syntax-rules, eval, environment, query },
  morekeywords=[3]{import, export,
    make, list, component, frequency, fwidth,
    harminv, output-efield-z },
  alsodigit=!\$\%&*+-./:<=>?@^_~,
  sensitive=true,
  morecomment=[l]{;},
  morecomment=[s]{\#|}{|\#},
  morestring=[b]",
  basicstyle=\small\ttfamily,
  keywordstyle=\bf\ttfamily\color[rgb]{0,.3,.7},
  commentstyle=\color[rgb]{0.133,0.545,0.133},
  stringstyle={\color[rgb]{0.75,0.49,0.07}},
  upquote=true,
  breaklines=true,
  breakatwhitespace=true,
  literate=*{`}{{`}}{1}
}

\definecolor{dkgreen}{rgb}{0,0.6,0}
\definecolor{gray}{rgb}{0.5,0.5,0.5}
\definecolor{mauve}{rgb}{0.58,0,0.82}

\lstset{ %
  language=Scheme,                % the language of the code
  basicstyle=\footnotesize,       % the size of the fonts that are used for the code
  numbers=left,                   % where to put the line-numbers
  numberstyle=\tiny\color{gray},  % the style that is used for the line-numbers
  stepnumber=1,                   % the step between two line-numbers. If it's 1, each line 
                                  % will be numbered
  numbersep=5pt,                  % how far the line-numbers are from the code
  backgroundcolor=\color{white},  % choose the background color. You must add \usepackage{color}
  showspaces=false,               % show spaces adding particular underscores
  showstringspaces=false,         % underline spaces within strings
  showtabs=false,                 % show tabs within strings adding particular underscores
  frame=single,                   % adds a frame around the code
  rulecolor=\color{black},        % if not set, the frame-color may be changed on line-breaks within not-black text (e.g. commens (green here))
  tabsize=4,                      % sets default tabsize to 2 spaces
  captionpos=b,                   % sets the caption-position to bottom
  breaklines=true,                % sets automatic line breaking
  breakatwhitespace=false,        % sets if automatic breaks should only happen at whitespace
  title=\lstname,                 % show the filename of files included with \lstinputlisting;
                                  % also try caption instead of title
  keywordstyle=\color{blue},          % keyword style
  commentstyle=\color{dkgreen},       % comment style
  stringstyle=\color{mauve},         % string literal style
  escapeinside={\%*}{*)},            % if you want to add a comment within your code
  morekeywords={*,...}               % if you want to add more keywords to the set
}
}

% -----------------------------------------------------------------------------
% Custom styling for quote boxes, use \epigraph

\definecolor{quotationcolour}{HTML}{F0F0F0}
\definecolor{quotationmarkcolour}{HTML}{1F3F81}

\newcommand{\epiline}{\hrule \vskip -.2em \hrule}
% Massively humongous opening quotation mark.
\newcommand{\hugequote}[1]{%
  \fontsize{42}{48}\selectfont \color{quotationmarkcolour} \textbf{#1}
  \vskip -.5em
}

\newcommand{\epigraph}[2]{%
  \smallskip
 \begin{center}
%  \hspace{1em}
  \colorbox{quotationcolour}{%
    \parbox{.8\textwidth}{%
    \epiline \vskip 1em 
    #1
    \begin{flushright}\textsc{#2}\end{flushright}
    \epiline
    }
  }
  \end{center}
  \medskip
}

\newcommand{\epigraphshort}[1]{%
  \smallskip
 \begin{center}
%  \hspace{1em}
  \colorbox{quotationcolour}{%
    \parbox{.8\textwidth}{%
    \epiline \vskip 1em 
     \begin{center}
    #1
    \end{center}

    \epiline
    }
  }
  \end{center}
  \medskip
}

% -----------------------------------------------------------------------------
% custom DRAFT MARK

\usepackage[]{datetime} 
\usepackage[allpages]{draftmark}

\iffinal{}{
\draftmarksetup{
  angle=0,scale=.34,
  color=red!45!blue!25,xcoord=117,ycoord=-129,
  mark={  \raggedright \mmddyyyydate (DRAFT)\\ \today} }
}

% -----------------------------------------------------------------------------
% Customize fncychap settings
\iffinal{}{
	\ChRuleWidth{1.5pt}
	\ChTitleLowerCase
	\ChTitleVar{\Huge \sc}
	}

% -----------------------------------------------------------------------------
% tighten the fncychap spacings, TeX code taken from
% http://tex.stackexchange.com/questions/13357/fncychap-package-reduce-vertical-gap-space-between-header-and-chapter-heading

\makeatletter
\renewcommand*{\@makechapterhead}[1]{%
  \iffinal{\setstretch{\@ssp}}{}
  
  \vspace*{\chaptertopspacing\p@}%
  {\parindent \z@ \raggedright \normalfont
    \ifnum \c@secnumdepth >\m@ne
      \if@mainmatter%%%%% Fix for frontmatter, mainmatter, and backmatter 040920
        \DOCH
      \fi
    \fi
    \interlinepenalty\@M
    \if@mainmatter%%%%% Fix for frontmatter, mainmatter, and backmatter 060424
      \DOTI{#1}%
    \else%
      \DOTIS{#1}%
    \fi
    \vspace*{\chaptertotextspacing\p@}
  }
  \iffinal{\setstretch{\@dsp}}{}  
}


 \renewcommand*{\@makeschapterhead}[1]{%
  \vspace*{\chaptertopspacingSTAR\p@}%
  {\parindent \z@ \centering%\raggedright
    \normalfont
    \interlinepenalty\@M
    \DOTIS{#1}
    \vskip \chaptertotextspacingSTAR\p@
  }}
 
\makeatother


% -----------------------------------------------------------------------------
% redefine abstract from drexel-thesis since the spacing is now wrong
\makeatletter
\renewenvironment{abstract}{%
  \listed@schapter{\abstractname}%
  \blanklines{-0}%
    \begin{center}
      \setstretch{\@ssp}%
      \@DUT@title\\
      \@DUT@author\\
      \ifdaring{%
        \ifnum\c@@DUT@advisors=\@ne%
        Advisor:
        \else%
        Advisors:
        \fi}{}
      \@DUT@advisor\\
    \end{center}
  \blanklines{4}%
  \setstretch{\@dsp}%
  \@nobreaktrue
  \@afterindentfalse
  \@afterheading
}{%
  \par\setstretch{\@ssp}%
}
\makeatother
